\newpage
\begin{center}
  \textbf{\large 4 Методические рекомендации}
\end{center}
\refstepcounter{chapter}
\addcontentsline{toc}{chapter}{4 Методические рекомендации}

\section{Рекомендации по использованию}
На основе опыта, полученного при проведении экспериментов с методом DVLO и предложенной модификацией 
можно выделить следующие рекомендации по использованию:
\begin{itemize}
  \item В качестве оборудования для получения данных стоит выбирать камеры и лидар высокого 
  разрешения для обеспечения возможности экстракции большего количества признаков 
  и учета их при регрессии положения системы в ходе работы алгоритма.
  \item Ключевую роль в работоспособности метода играет синхронизация данных. Наличие
  смещения по времени между показаниями датчиков может повлечь искажение проекции лидарных точек в 
  систему координат камеры, что усложнит процесс обучения и приведет к снижению качества одометрии.
  \item На качество проекции при экстракции признаков и локальном объединении влияет также взаимная 
  калибровка камеры и лидара. В связи с этим рекомендуется тщательно верифицировать параметры системы сенсоров.
  \item На начальном этапе обучения важно обеспечить равный вклад трансляционной ошибки и ошибки ориентации в 
  функцию потерь, для этого значения адаптивных коэффициентов $k_x$ и $k_q$ выбираются исходя из условия совпадения 
  масштаба  данных значений. Для различных данных может потребоваться пересмотр предлагаемых в оригинальной статье начальных условий.
\end{itemize}

\section{Потенциальные улучшения}
\label{sec:future_work}

Ниже представлены наиболее перспективные расширения,
которые могут существенным образом повысить точность и робастность одометрии. 
Кроме того стоит отметить, что использование полностью дифференцируемого алгоритма 
позволяет значительно расширять ее функционал путем добавления параллельных блоков, 
использующих признаки из основной модели.

\subsection{Учёт временного контекста посредством трансформера}
\label{subsec:temporal_transformer}

В текущей реализации DVLO / Stereo‑DVLO каждая пара
последовательных кадров обрабатывается независимо.
Это не позволяет модели использовать дополнительную динамическую информацию, 
что могло бы позволить обрабатывать более сложные профили движения.
Возможным решением является использование 
временного трансформера (англ. Temporal Transformer), принимающего в окно несколько 
последних глобально объединенных признаков ${F}_{G}$.
Механизм внимания может обеспечить адаптивный учет вкладов прошлых кадров,
выделяя длительные закономерности в данных и подавляя шум.

\subsection{Использование оценки глубины}
\label{subsec:stereo_depth}

Помимо прямой интеграции второй камеры в оригинальную структуру DVLO 
для учета текстурной информации, предлагается первично оценивать глубину при 
помощи компактной нейронной сети (например H-Net~\cite{huang2022h}). Далее
при локальном объединении помимо цветовых признаков каждой точки используется 
также признак глубины. Такая схема позволяет сохранить парадигму полностью нейросетевого 
подхода одометрии, при этом также позволяет расширить его применимость за счет повышения 
точности и функциональности (оценки глубины можно использовать независимо для параллельных задач).


